% Figure: a hinged trapdoor held open by a cable. The door is a uniform slab
% hinged at one edge, propped open by a cable to a wall point above. FBD shows
% weight at the centroid, hinge reaction, cable tension. Numbers match the
% trapdoor worked example.
\begin{figure}[htbp]
\centering
\begin{tikzpicture}[
  force/.style={draw=engblue, -{Stealth}, very thick},
  rxn/.style={draw=engred, -{Stealth}, very thick},
  cable/.style={draw=engamber, line width=1.4pt},
  scale=1.0,
]
% wall (vertical)
\draw[draw=engblack, very thick] (0,-0.4) -- (0,3.4);
\foreach \y in {-0.3,0.1,...,3.3} {\draw[draw=enggray, thin] (0,\y) -- (-0.26,\y-0.2);}
% hinge at H at wall base
\coordinate (H) at (0,0);
\draw[draw=engblack, fill=enggray!20, thick] (H) circle (0.12);
\node[font=\scriptsize, anchor=north east] at (H) {$H$};
% door raised at angle phi above horizontal, length Ld
\def\phi{35}
\coordinate (T) at ({3.4*cos(\phi)},{3.4*sin(\phi)});
\draw[draw=engblack, line width=2.4pt] (H) -- (T);
% centroid
\coordinate (G) at ({1.7*cos(\phi)},{1.7*sin(\phi)});
\fill[engblack] (G) circle (0.06);
% cable from wall point C (up the wall) to door tip T
\coordinate (C) at (0,3.0);
\draw[cable] (C) -- (T);
\draw[draw=engblack, fill=enggray!20, thick] (C) circle (0.08);
\node[font=\scriptsize, anchor=east] at (C) {$C$};
% weight at centroid
\draw[force] (G) -- ++(0,-1.5) node[anchor=north, font=\small, engblue] {$W$};
% cable tension along door tip toward C
\draw[rxn] (T) -- ($(T)!0.42!(C)$) node[anchor=south west, font=\small, engred] {$T$};
% hinge reaction (schematic)
\draw[rxn] (H) -- ++(1.1,0.7) node[anchor=south west, font=\small, engred] {$R_H$};
% angle phi at hinge
\draw[draw=enggray, thin] (0.9,0) arc (0:\phi:0.9);
\node[font=\scriptsize, enggray] at (1.05,0.32) {$\phi$};
\node[font=\scriptsize, anchor=south west] at (T) {tip};
\end{tikzpicture}
\caption[A hinged trapdoor held open by a cable]{The chosen body is the
trapdoor, a uniform slab hinged at $H$ and propped open at angle $\phi$ above
horizontal by a cable running from the door tip to the wall point $C$. Three
external forces act: the weight $W$ down at the centroid, the cable tension
$T$ along the tip-to-$C$ line, and the hinge reaction $R_H$. Summing moments
about the hinge kills the hinge reaction and leaves one equation for the cable
tension. The weight's moment arm about $H$ shrinks with the horizontal
projection of the door, so a door raised nearer the vertical needs less cable
tension to hold, while a nearly horizontal door demands the most. The cable's
own geometry sets how much of its tension does useful lifting: a cable pulling
nearly along the door has a small moment arm and must carry a large tension to
supply the holding moment.}
\label{fig:vol03ch01:trapdoor}
\end{figure}
